% Folder 142, batch 02 — archivists' pages 21–40.
% Pass: Opus 5.5 (claude-opus-5-5), 2026-10-02 — first pass, unchecked against the pages by a human.
%
% Page 21 is a cover in his hand; its words are the section heading. Page 22
% carries only a pencilled page count and is skipped. Pages 23–40 are one run
% paginated by him 5–22 (formulas (1)–(79)); it continues past page 40.
\documentclass[11pt,a4paper]{article}
\input{../preamble/grothendieck}

\folder{142}
\batch{2}
\pages{21}{40}
\foldertitle{Action de Galois sur Teichmüller : notes manuscrites (s.d.).}
\dating{[à partir de 1978]}
\watermark{Édition de démonstration}

\begin{document}

\section{Opérations de $\Gamma_{\overline{\mathbb{Q}}/\mathbb{Q}}$ sur $\pi_1 M_{0,3}$}
\note{titre de la main de Grothendieck, sur le feuillet de couverture (p.~21), d'ailleurs blanc ; un symbole biffé \ill{} précède $\pi_1$. Le texte qui suit commence à sa p.~5, mais la numérotation des formules y repart de (1) : les pp.~1–4 de l'auteur ne sont pas dans ce lot.}

\page{23}
\note{p.~5 de l'auteur}
\drawing{sphère vue en perspective ; sur l'équateur, trois points marqués (ronds pleins), chacun entouré d'un petit cercle portant deux sommets $R_i^{\omega}$, $R_i^{\omega'}$ ($i = 0, 1, \infty$), reliés par les arcs $a_i^{\omega}$, $a_i^{\omega'}$ ; les sommets voisins de deux points différents sont reliés par les flèches $b_0^{\omega}$, $b_1^{\omega}$ (lecture douteuse), $b_\infty^{\omega}$. Légende : $\mathrm{I}_J$.}
\marginal{NB. $\tau(R_i^{\omega}) = R_i^{\omega'}$, $\tau(a_i^{\omega}) = (a_i^{\omega'})^{-1}$, $\tau(b_i^{\omega}) = b_i^{\omega'}$}
\note{dans cette note marginale, écrite en biais, les exposants $\omega'$ sont peu lisibles.}

$\widetilde{\mathfrak{S}}_J$ permute transitivement les six flèches $a_i^{\omega}$ et les six flèches $b_i^{\omega}$ : \struck{\ill{}}
\[
(1)\quad
\begin{cases}
a_i^{\omega} : R_i^{\omega} \to R_i^{\omega'} \\
b_i^{\omega} : R^{\omega}_{\omega i} \to R^{\omega'}_{\omega' i}
\end{cases}
\]
\[
(2)\quad
\begin{cases}
\text{a)}\ \ b_i^{\omega} b_i^{\omega'} = 1 \\
\text{b)}\ \ a_i^{\omega} b^{\omega}_{\omega' i} a_i^{\omega'} b^{\omega}_{\omega i} a^{\omega'}_{\omega i} b_i^{\omega} = 1
\end{cases}
\]
(trois relations \uncertain{effectives}) ; (deux relations effectives, une pour chaque $\omega \in \underline{\omega}$)
\note{les indices et exposants de (2\,b) sont lus avec doute.}
\[
(3)\quad a_i^{\omega'} a_i^{\omega} = \lambda_i^{\omega} = \varphi_i^{\omega}(l_0)
\]
\[
(4)\quad \varphi_i^{\omega} : \Pi_{0,3} = \pi_1(M_{0,3}\,;R_0^{+}) \hookrightarrow \pi_1(\Sigma^{*}, R_i^{\omega})
\]
\[
\Pi_{0,3} = \{\, l_0, l_1, l_\infty \mid l_\infty l_1 l_0 = 1 \,\}
\subset \widetilde{\Pi}_{0,3} = \pi_1(M_{0,3}, \mathfrak{S}_3\,; R_0^{+}),
\]
\[
l_0 \mapsto a_i^{\omega'} a_i^{\omega}, \qquad
l_1 \mapsto (b^{\omega}_{\omega' i})^{-1} a^{\omega}_{\omega i} a^{\omega'}_{\omega i} b_i^{\omega}
\]
\note{les indices de l'image de $l_1$ sont incertains ; un trait sous l'expression porte « \ill{} ».}
Les $\varphi_i^{\omega}$ sont permutés entre eux de façon simplement transitive.
\[
(5)\quad \tau\varphi_i^{\omega}(g) = \varphi_i^{\omega}(\tau_\infty(g)) \quad \text{où } \tau_\infty \in \widetilde{\Pi}_{0,3},
\]
\[
\tau_\infty(l_0) = l_0^{-1}, \quad \tau_\infty(l_1) = l_1^{-1}, \quad
\tau_\infty(l_\infty) = l_0 l_1 = \mathrm{int}(l_0)\, l_\infty^{-1}
\]
\[
(6)\quad
\begin{cases}
a_i^{\omega}(\varphi_i^{\omega}(g)) = \varphi_i^{\omega'}(\varepsilon_0(g)) \\
b_i^{\omega}(\varphi^{\omega}_{\omega i}(g)) = \varphi^{\omega'}_{\omega' i}(\sigma_\infty(g))
\end{cases}
\qquad \varepsilon_0^2 = l_0
\]
\[
\sigma_\infty(l_0) = l_1, \quad \sigma_\infty(l_1) = l_0, \quad
\sigma_\infty(l_\infty) = (l_0 l_1)^{-1} = l_1^{-1} l_0^{-1} = \mathrm{int}(l_0)\, l_\infty
\]
\struck{\ill{}}

On introduit
\[
(7)\quad \tilde b_i^{\omega} = a^{\omega'}_{\omega' i} b_i^{\omega} : R^{\omega}_{\omega i} \to R^{\omega}_{\omega' i} ;
\]
(2\,b) s'écrit
\[
(8)\quad \tilde b^{\omega}_{\omega^2 i}\, \tilde b^{\omega}_{\omega i}\, \tilde b_i^{\omega} = 1
\]
\marginal{NB. L'un des $\tilde b_i^{\omega}$ n'est pas stable sous $\tau$, il faudrait donc \uncertain{aussi} prendre les $\tau \tilde b_i^{\omega}$}
On aura
\[
(9)\quad \tilde b_i^{\omega}(\varphi^{\omega}_{\omega i}(g)) = \varphi^{\omega}_{\omega' i}(\rho^{-1}(g))
\qquad \rho^{-1} = \varepsilon_0 \sigma_\infty,
\quad
\begin{cases}
\rho(l_0) = l_1 \\
\rho(l_1) = l_\infty \\
\rho(l_\infty) = l_0
\end{cases}
\]
Les relations 2\,b) deviennent
\[
(\text{\struck{\ill{}}})\quad \tilde b_{\omega^2 i}\, \tilde b_{\omega i}\, \tilde b_i = 1
\]

\page{24}
\note{p.~6 de l'auteur}
\uncertain{Soit} $u$ \uncertain{un} automorphisme \add{du groupoïde $\widehat{\mathrm{I}}_J$} qui commute à $\widetilde{\mathfrak{S}}_J$, et fixe les pts, \struck{il est déterminé} et qui \uncertain{conserve} \ill{} les sous-groupoïdes $\widehat{\mathrm{I}}_{J,i}$, où $\widehat{\mathrm{I}}_{J,i}$ est engendré par $R_i^{\omega}$, $a_i^{\omega}$, $\omega \in \underline{\omega}$. La restriction à $\mathrm{I}_{J,i}$ \struck{définie} \uncertain{est} donnée par
\[
u(a_i^{\omega}) = a_i^{\omega} \varphi_i^{\omega}(l_0^{\gamma}) = \varphi_i^{\omega'}(l_0^{\gamma})\, a_i^{\omega}
\qquad \text{($\varphi_i^{\omega}(l_0^{\gamma}) = (\lambda_i^{\omega})^{\gamma}$)}
\]
$l_0^{\gamma}$ \uncertain{indépendant} de $i, \omega$, $\gamma \in \hat{\mathbb{Z}}$ ; on aura
\[
u(a_i^{\omega'} a_i^{\omega}) = \varphi_i^{\omega}(l_0^{2\gamma+1}), \qquad
u(\underbrace{\varphi_i^{\omega}(l_0)}_{\lambda_i^{\omega}}) = \varphi_i^{\omega}(l_0^{p}) = (\lambda_i^{\omega})^{p},
\qquad \gamma = \frac{p-1}{2}
\]
\note{dans $l_0^{2\gamma+1}$ un caractère surchargé précède le $+1$.}
\[
\begin{cases}
u(a_i^{\omega}) = a_i^{\omega} \varphi_i^{\omega}(l_0^{\gamma}) \\
u(b_i^{\omega}) = b_i^{\omega} \varphi^{\omega}_{\omega i}(\beta)
\quad (= \varphi^{\omega'}_{\omega' i}(\sigma_\infty \beta)\, b_i^{\omega})
\end{cases}
\qquad \gamma = \frac{p-1}{2}, \quad p \in \hat{\mathbb{Z}}^{*} \text{ multiplicateur}, \quad \beta \in \widehat{\Pi}_{0,3}
\]
\note{sous la seconde ligne, une variante biffée et illisible.}
Ici $(p, \beta)$ détermine $u$ (en fait, $\beta$ détermine $p$…)
\[
u(\tilde b_i^{\omega}) = \tilde b_i^{\omega} \varphi^{\omega}_{\omega i}(\underbrace{l_1^{\gamma} \beta}_{\sigma_\infty(\lambda) \beta})
= \varphi^{\omega'}_{\omega' i}(l_0^{\gamma} \rho^{-1}(\beta))\, \tilde b_i^{\omega}
\]
Les équations sur $(p, \beta)$ \uncertain{expriment} 2 \ill{}
\[
(12)\quad
\begin{cases}
\beta\, \sigma_\infty(\beta) = 1 \\
\alpha\, \rho^{-1}(\alpha)\, \rho^{-2}(\alpha) = 1
\quad \text{où } \text{\struck{$\alpha = l_1^{-\gamma} \beta$}} \\
\alpha = \underbrace{l_1^{\gamma}}_{\sigma_\infty(\lambda)} \beta
\end{cases}
\]
\marginal{NB. Si \uncertain{on remplace} $\beta$ par $\beta^{-1}$, $\alpha$ par $\alpha^{-1}$, on retrouve les équations habituelles ($\alpha^{-1} \rho(\alpha^{-1}) \rho^{2}(\alpha^{-1}) = 1$) — cf plus \uncertain{loin}}
Notons que si $u$ est défini par $(p, \beta)$, alors $\tau(u)$ \add{$= \tau u \tau^{-1}$} est défini par $(p, \beta' = \tau_\infty(\beta))$. Il faudrait pouvoir \uncertain{vérifier} directement, par un calcul dans $\Pi_{0,3}$, que les équations (12) sur $\beta$ impliquent les \uncertain{mêmes} relations pour $\beta' = \tau_\infty(\beta)$ (avec \uncertain{le même} $\gamma$).
\marginal{(en haut, en biais) \uncertain{Les deux automorphismes} \ill{} \uncertain{devrait plus} \ill{} \uncertain{on ne sait pas} \ill{} $(p, \beta)$, \ill{} $p \in \hat{\mathbb{Z}}$ \ill{} $p \in \hat{\mathbb{Z}}^{*}$, \ill{} $2\gamma+1$ \ill{} $\gamma \in \hat{\mathbb{Z}}$}
\marginal{(au milieu, en biais) Faire les calculs \uncertain{d'abord} \ill{} $u(a_i^{\omega}) = a_i^{\omega} \varphi_i^{\omega}(\lambda)$ avec $\lambda \in \widehat{\Pi}_{0,3}$ \ill{} de la forme $l_0^{\gamma}$, $\gamma \in \hat{\mathbb{Z}}$…}
\marginal{(en bas, en biais) Si on change le \uncertain{définition} de $\beta$, \ill{} $\beta$ devient $\sigma_\infty(\beta)$ \ill{} les deux premières \ill{} $\alpha = l_1^{\gamma} \sigma_\infty(\beta)$ \ill{}}
\note{les trois notes marginales de cette page, écrites en biais et très serrées, ne se lisent que par fragments.}

\page{25}
\note{p.~7 de l'auteur}
\drawing{même sphère que p.~23, à laquelle s'ajoutent sur l'équateur trois nouveaux sommets $Q_0$, $Q_1$, $Q_\infty$ et des arcs $c_i^{\pm}$ qui les relient aux sommets $R^{\pm}$ voisins. Légende : $\mathrm{II}_J$.}
On rajoute les sommets $Q_i$ ($i \in J$) et les arcs
\[
Q_i \xrightarrow{\;c_i^{\omega}\;} R^{\omega'}_{\omega' i}
\]
permutés \struck{transitivement} \add{de façon simple\uncertain{ment}} \uncertain{transitive} par $\widetilde{\mathfrak{S}}_J$. On a maintenant
\[
(13)\quad b_i^{\omega} = c_i^{\omega} (c_i^{\omega'})^{-1}
\]
(donc les $b_i^{\omega}$ deviennent inutiles comme générateurs d'un groupoïde). (\add{Donc} les relations (2) \struck{deviennent} \ill{}, la partie (a) devient automatique, pour écrire commodément la partie (b), on introduit les
\[
(14)\quad \tilde a_i^{\omega} = (c^{\omega}_{\omega i})^{-1} a_i^{\omega} c^{\omega'}_{\omega' i} : Q_{\omega i} \to Q_{\omega' i}
\]
\note{la première écriture de (14), surchargée, est remplacée par celle, interlinéaire, qui est donnée ici.}
\drawing{à droite : petit schéma de $\tilde a_i^{\omega}$, chemin de $Q_{\omega i}$ à $Q_{\omega' i}$ passant par $(c^{\omega}_{\omega i})^{-1}$, l'arc $a_i^{\omega}$ entre $R_i^{\omega}$ et $R_i^{\omega'}$, puis $c^{\omega'}_{\omega' i}$ ; une première version, au-dessous, est rayée. Dans la marge gauche, un croquis analogue, en partie rayé.}
qui sont également permutés de façon simpl. transitive par $\widetilde{\mathfrak{S}}_J$, et on a
\[
(15)\quad \tau(\tilde a_i^{\omega}) = (\tilde a_i^{\omega'})^{-1}
\]
Les relations (2\,b) deviennent
\[
(16)\quad \tilde a^{\omega}_{\omega^2 i}\, \tilde a^{\omega}_{\omega i}\, \tilde a_i^{\omega} = 1
\]
(deux relations ess\supplied{entiellemen}t distinctes, qui sont permutées par $\tau$)

Pour repérer les \add{$\pi(Q_i) =$} $\pi_1(\Sigma_J(\mathbb{C}), Q_i)$ de façon commode, on utilise les
\[
(17)\quad \psi_i^{\omega} : \Pi_{0,3} \xrightarrow{\ \sim\ } \pi(Q_i)
\]

\begin{tikzcd}[column sep=small]
\Pi_{0,3} \arrow[rr, "\sim"] \arrow[dr, "{\varphi^{\omega'}_{\omega' i}\,\wr}"'] & & \pi(Q_i) \\
& \pi(R^{\omega'}_{\omega' i}) \arrow[ur, "{(c_i^{\omega})^{-1},\,\sim}"'] &
\end{tikzcd}

\marginal{NB. $\tau\psi_i^{\omega}(g) = \psi_i^{\omega}(\tau_\infty g)$}
\note{dans ce NB l'exposant du second membre est peu sûr.}
on aura
\[
(18)\quad \psi_i^{\omega}(g) = \psi_i^{\omega'}(\sigma_\infty(g))
\]
\[
(19)\quad
\begin{aligned}
\tilde a_i^{\omega}(\psi^{\omega'}_{\omega i}(g)) &= \psi^{\omega}_{\omega' i}(\varepsilon_0(g)) = \psi^{\omega'}_{\omega' i}(\sigma_\infty \varepsilon_0(g)) \\
\tilde a_i^{\omega}(\psi^{\omega}_{\omega i}(g)) &= \psi^{\omega}_{\omega' i}(\underbrace{(\varepsilon_0 \sigma_\infty)}_{\rho^{-1}}(g)) = \psi^{\omega'}_{\omega' i}(\sigma_\infty \varepsilon_0 \sigma_\infty(g))
\end{aligned}
\]
\note{dans (19), plusieurs exposants $\omega$ sont surchargés d'un astérisque ; au-dessus de la première ligne, « $\tau_\infty(\rho)$ » précédé d'un mot biffé. La seconde ligne, encadrée, semble corriger la première.}

\page{26}
\note{p.~8 de l'auteur}
On prend un automorphisme \add{$u$} des groupoïdes $\widehat{\mathrm{II}}_J$ qui commute à $\widetilde{\mathfrak{S}}_J$ et qui fixe les sommets, et qui conserve plus \uncertain{exactement} \uncertain{bien stables} les sous-groupoïdes $\widehat{\mathrm{II}}_{J,i}$ ($i \in J$), \struck{engendrés de} $\mathrm{II}_{J,i}$ engendré par les \add{deux} sommets $R_i^{\omega}$ ($\omega \in \underline{\omega}$) et les deux flèches $a_i^{\omega}$ ($\omega \in \underline{\omega}$). \struck{\ill{}} L'automorphisme $u$ est connu quand on connaît les $u(a_i^{\omega})$, qui sont de la forme
\[
(20)\quad u(a_i^{\omega}) = a_i^{\omega} \varphi_i^{\omega}(l_0^{\gamma})
\qquad \gamma = \frac{p-1}{2} \in \hat{\mathbb{Z}}, \quad p \in \hat{\mathbb{Z}}^{*}
\]
\marginal{NB : $\lambda \in \widehat{\Pi}_{0,3}$, \uncertain{pas le peine d'}exiger de prime abord que $\lambda$ \uncertain{soit} de la forme $l_0^{\gamma}$…}
\emph{et} les
\[
(21)\quad u(c_i^{\omega}) = \varphi^{\omega'}_{\omega' i}(\eta)\, c_i^{\omega} = c_i^{\omega} \psi_i^{\omega}(\eta)
\qquad \eta \in \widehat{\Pi}_{0,3} \text{ indépendant de } i, \omega
\]
On aura alors (comparer (11))
\[
(22)\quad u(b_i^{\omega}) = b_i^{\omega} \varphi^{\omega}_{\omega i}(\beta)
\]
avec
\[
(23)\quad \beta = \sigma_\infty(\eta)\, \eta^{-1}
\]
\marginal{NB. Devient plus joli si $\eta$ est remplacé par \uncertain{son} $\sigma_\infty(\eta)$, ou $\beta$ par $\beta^{-1}$}
\marginal{(en biais, sous (23)) $\eta' = \tau(\eta)$, $\beta' = \tau_\infty(\beta)$ \uncertain{?}}
Les relations à satisfaire sont celles déjà écrites sur $p$ (ou $\gamma$) et $\beta$, \uncertain{savoir} (12) — mais la première relation $\beta\sigma_\infty(\beta) = 1$ \uncertain{est} conséquence de (23), il reste la deuxième, qui s'écrit donc
\[
(24)\quad \text{\struck{$l_1^{-\gamma}\sigma_\infty(\eta)\eta^{-1}$}}\ \ \alpha \rho^{-1}(\alpha) \rho^{-2}(\alpha) = 1
\quad \text{avec } \alpha = \underbrace{l_1^{\gamma}}_{\sigma_\infty(\lambda)} \underbrace{\sigma_\infty(\eta)\eta^{-1}}_{\beta}.
\]
\marginal{NB. Si $u$ est remplacé par $\tau u \tau^{-1}$, alors $\eta$ \uncertain{est} remplacé par $\tau_\infty(\eta)$, $\beta$ par $\varepsilon_0^{-1}(\tau_\infty(\lambda^{-1}))$ \ill{}}
\note{note marginale en biais, lecture très incertaine.}
On aura ici
\[
(25)\quad u(\tilde a_i^{\omega}) = \psi^{\omega^*}_{\omega i}(\underbrace{\eta^{-1} l_0^{\gamma} \varepsilon_0(\eta)}_{\tilde\alpha = \eta^{-1}\varepsilon_0(\lambda\eta)})\, \tilde a_i^{\omega}
= \tilde a_i^{\omega}\, \psi^{\omega'}_{\omega' i}(\underbrace{\varepsilon_0^{-1}(\eta^{-1}) l_0^{\gamma} \eta}_{\varepsilon_0^{-1}(\tilde\alpha)})
\]
\note{sous $\tilde\alpha$ : « valable pour tout $\lambda$ ». Dans la marge gauche : $\tilde\alpha = \eta^{-1} l_0^{\gamma} \varepsilon_0(\eta)$.}
et
\[
(\text{25\,bis})\quad \alpha = \rho(\eta \tilde\alpha \rho^{-1}(\eta)^{-1})
\]
valable pour \uncertain{$\lambda$} quelconque, pas \uncertain{nécessairement} de la forme $l_0^{\gamma}$. \struck{\ill{}} Dans l'équation (24) on \uncertain{a} (\uncertain{invariante} par substitutions $\alpha \mapsto \rho^{-1}(\alpha)$ (ou $\alpha \mapsto \rho\alpha = \rho^{-2}(\alpha)$) et $\alpha \mapsto x\alpha\rho^{-1}(x)^{-1}$,

\page{27}
\note{p.~9 de l'auteur}
équivalente à l'équation \struck{identique} \struck{\ill{}} \add{\uncertain{analogue}} en $\tilde\alpha$
\[
(26)\quad \tilde\alpha\, \rho^{-1}(\tilde\alpha)\, \rho^{-2}(\tilde\alpha) = 1, \qquad
\boxed{(27)\quad \tilde\alpha = \eta^{-1} l_0^{\gamma} \varepsilon_0(\eta)}
\]
\note{le second $\rho$ de (26) porte un exposant peu net, lu $-2$ d'après (24).}

Considérons maintenant le groupoïde \add{$\mathrm{III}_J$} induit par le groupoïde précédent \add{$\mathrm{II}_J$} sur l'un des trois sommets $Q_i$, c'est-à-dire le groupoïde engendré par le graphe à trois sommets $Q_i$ ($i \in J$) et six arcs orientés $\tilde a_i^{\omega}$ et leurs opposés, avec les seules relations (16) (une des 6 relations stable par $\widetilde{\mathfrak{S}}_J$, et se réduisant en fait à deux relations effectives…). Le groupe $\{1, \tau\}$ opère dans le graphe \uncertain{avec} relations, par
\[
(27)\quad \tau Q_i = Q_i, \qquad \tau \tilde a_i^{\omega} = (\tilde a_i^{\omega'})^{-1}
\]
\note{le numéro (27) sert ici une seconde fois, comme sur la page.}
\drawing{à gauche : sphère portant les trois sommets $Q_0$, $Q_1$, $Q_\infty$ et les six arcs $\tilde a_i^{\pm}$ qui les joignent deux à deux. Légende : $\mathrm{III}_J$.}
La donnée d'un \struck{automorphisme} \add{\ill{}} des groupoïdes $\widehat{\mathrm{III}}_J$, fixant les sommets et commutant à $\widetilde{\mathfrak{S}}_J$, \uncertain{équivaut} à celle d'un $\tilde\alpha \in \widehat{\Pi}_{0,3}$, \struck{\ill{}} \struck{l'équation \ill{}} satisfaisant (26). On aura, en introduisant
\[
(28)\quad \tilde\lambda_i^{\omega} = \tilde a_i^{\omega'} \tilde a_i^{\omega} \in \pi(Q_{\omega' i})
\qquad (\lambda_i^{\omega} = a_i^{\omega'} a_i^{\omega})
\]
\drawing{à droite : $Q_{\omega i}$ et $Q_{\omega' i}$ reliés par les deux arcs $\tilde a_i^{\omega}$ (en haut) et $\tilde a_i^{\omega'}$ (en bas) ; au centre, en pointillé, le petit cercle $R_i^{\omega}$, $R_i^{\omega'}$ et l'arc $a_i^{\omega}$.}
\marginal{(contourner) \uncertain{\ill{} des $\tilde a_i^{\omega}$ \ill{} groupoïde $\mathrm{II}_J$ !} \ill{}}
\[
(29)\quad u(\tilde\lambda_i^{\omega}) = \psi^{\omega'}_{\omega' i}(\tilde\alpha\, \varepsilon_0(\tilde\alpha))\, \tilde\lambda_i^{\omega}
= \psi^{\omega'}_{\omega' i}(\tilde\alpha\, \varepsilon_0(\tilde\alpha)\, l_0)
\]
Dans le cas où $\tilde\alpha$ provient de $\eta, \gamma$ par (27), on a
\[
(30)\quad \tilde\alpha\, \varepsilon_0(\tilde\alpha)\, l_0 = \eta^{-1} l_0^{p} \eta \qquad (p = 2\gamma+1) ;
\]
\marginal{(en biais, à gauche) $\psi$ \uncertain{serait} \uncertain{raisonnable} \ill{} les $\tilde a_i^{\omega}$ \ill{} dans $\mathrm{II}_J$ \ill{} — * \ill{}}
\note{le numéro (30) est écrit dans la marge, en bas à gauche ; les notes marginales de cette page sont écrites en biais et se lisent mal.}

\page{28}
\note{p.~10 de l'auteur}
ou, en d'autres termes, faisant opérer $u$ sur $\pi(Q_{\omega' i})$, identifié à $\widehat{\Pi}_{0,3}$ via $\psi^{\omega'}_{\omega' i}$, $\tilde\lambda_i^{\omega}$ identifié à $l_0$, on trouve que $u(l_0)$ est conjugué de $l_0^{p}$, \struck{\ill{}} la conjugaison étant effectuée par $\eta^{-1}$ (par la formule (30)).
\marginal{(en biais) NB. Le \uncertain{sous-groupoïde} $\widehat{\mathrm{III}}_{J,i}$ (\ill{}) \uncertain{engendré} \ill{} \uncertain{est stable par} $u$ \ill{} $\mathfrak{S}_J \times \{1, \tau\}$, \ill{}}

Le passage de $u$ à $\tau(u)$ s'exprime par \struck{\uncertain{l'on rajoute maintenant les points} $P_\omega$ \ill{} ($\omega \in \underline{\omega}$)} la passage de $\tilde\alpha$ à
\[
(31)\quad \tilde\alpha' = \varepsilon_0 \tau_\infty(\tilde\alpha^{-1}) = \underbrace{\rho^{-1} \tilde\sigma_\infty}_{\tilde\sigma_0}(\tilde\alpha^{-1})
\]
où $\tilde\sigma_\infty = \sigma_\infty \tau_\infty = \tau_\infty \sigma_\infty$ est dans l'image de l'homomorphisme canonique $\widetilde{\mathfrak{S}}_3 \to \widetilde{\Pi}_{0,3}$, \struck{\ill{}} \uncertain{$\simeq$} $\widetilde{\mathfrak{S}}_3$ stabilisateur de $P_{+}$ dans $\mathfrak{S}_3 \times \{1, \tau\}$
\[
\tilde\sigma_i(l_i) = l_i^{-1}, \qquad
\tilde\sigma_i(l_j) = l_k^{-1} \quad \text{si } \{0, 1, \infty\} = \{i, j, k\}
\]

\textbf{NB.} Le groupe $\widetilde{\mathfrak{S}}_3$ opère sur $\Pi_{0,3}$ et contient $\widetilde{\mathfrak{S}}_3^{+}$ \struck{\ill{}} $\{1, \rho^{-1}, \rho^{-2}\}$ comme sous-groupe d'indice 2, \struck{\ill{} $\widetilde{\mathfrak{S}}_3$ \ill{}} $\widetilde{\mathfrak{S}}_3 / \widetilde{\mathfrak{S}}_3^{+}$ opère sur $\widetilde{\mathfrak{S}}_3^{+}$ en échangeant $\rho$ et $\rho^{-1}$, il s'ensuit que $\widetilde{\mathfrak{S}}_3$ opère sur l'ensemble des solutions $x$ de l'équation
\[
x\, \rho^{-1}(x)\, \rho^{-2}(x) = 1
\]
par
\[
(32)\quad \Theta_g x = g \cdot x^{\chi(g)}
\]
\note{l'exposant de (32), surchargé, est lu $\chi(g)$ sans certitude.}
et on a ici
\[
(33)\quad \tilde\alpha' = \Theta_{\tilde\sigma_0}(\tilde\alpha).
\]

\page{29}
\note{p.~11 de l'auteur}
\drawing{sphère : pôles $P_{+}$ (en haut) et $P_{-}$ (en bas) ; sur l'équateur, les sommets $Q_0$, $Q_1$, $Q_\infty$ et les trois points marqués (ronds pleins) ; les arcs $\tilde a_i^{\pm}$ le long de l'équateur, et les méridiens $d_i^{+}$ (de $P_{+}$ à $Q_i$) et $d_i^{-}$ (de $P_{-}$ à $Q_i$), en trait plein ou en pointillé selon qu'ils sont vus ou cachés.}
\marginal{NB. Les flèches $\tilde a_i^{\omega}$ sont ici superflues, cf (37) et commentaire p.~19.}
\note{la « p.~19 » est la pagination de l'auteur : p.~37 des archivistes.}
On rajoute les \add{deux} sommets $P^{\omega}$ ($\omega \in \underline{\omega}$), et les flèches
\[
d_i^{\omega} : P^{\omega} \to Q_i
\]
permutés transitivement par $\widetilde{\mathfrak{S}}_J$, avec
\[
(34)\quad \tau(P^{\omega}) = P^{\omega'}, \qquad \tau(d_i^{\omega}) = d_i^{\omega'}
\]
\note{dans (34) et dans la définition de $d_i^{\omega}$, l'indice de $P$ est surchargé ; on lit $P^{\omega}$.}

$\mathrm{IV}_J$\quad Le groupoïde $\mathrm{III}_J$ est équivalent au groupoïde $\mathrm{IV}_J$ \add{défini par le graphe \ill{}} \struck{sur} les sommets $Q_i$ ($i \in \omega$) et $P^{\omega}$ ($\omega \in \underline{\omega}$), les arcs $\tilde a_i^{\omega}$ et $d_i^{\omega}$ et les \add{arcs} opposés, et les six relations
\[
(35)\quad (d^{\omega}_{\omega' i})^{-1} \tilde a_i^{\omega} d_i^{\omega} = 1
\]
(qui impliquent d'ailleurs les relations correspondantes (16) pour $\mathrm{III}_J$). La donnée d'un \struck{automorphisme} \add{automorphisme} de $\widehat{\mathrm{IV}}_J$ qui fixe les pts sommets (pour simplifier — il faudrait regarder aussi ceux qui permutent les sommets \add{$P_\omega$}, en fixant les $Q_i$), \struck{équivalente} et qui commutent à $\widetilde{\mathfrak{S}}_J$, équivaut à la donnée des
\[
(36)\quad
\begin{cases}
u(\tilde a_i^{\omega}) = \psi^{\omega}_{\omega i}(\tilde\alpha)\, \tilde a_i^{\omega} & \tilde\alpha \in \widehat{\Pi}_{0,3} \ \text{(indép. de } i, \omega) \\
u(d_i^{\omega}) = \psi_i^{\omega}(\zeta)\, d_i^{\omega} \quad (= d_i^{\omega} \tilde\varphi_i^{\omega}(\zeta)) & \zeta \in \widehat{\Pi}_{0,3}\ \ldots
\end{cases}
\]
\note{la parenthèse de la seconde ligne de (36) est d'une lecture douteuse.}
de $\tilde\alpha, \zeta \in \widehat{\Pi}_{0,3}$, avec la relation déduite de (35), on trouve \struck{les} équations
\[
(37)\quad \tilde\alpha = \zeta \cdot \rho^{-1}(\zeta^{-1})
\]
i.e.\ $\tilde\alpha$ est connu, quand on connaît $\zeta \in \widehat{\Pi}_{0,3}$.
\marginal{(en biais) il n'y a donc \uncertain{plus de} conditions \uncertain{sur} $\zeta$ \uncertain{ici}.}

\page{30}
\note{p.~12 de l'auteur}
qui implique bien sûr (26).
\marginal{NB. Le passage de $u$ à $\tau(u)$ s'exprime par $\zeta \mapsto \sigma_\infty \tau_\infty(\zeta) = \tilde\sigma_\infty(\zeta)$}
\marginal{(en biais, à gauche) \uncertain{revenir à} l'introduction \ill{}}

Je \struck{vais} considère les isomorphismes

\begin{tikzcd}[column sep=small]
\Pi_{0,3} \arrow[rr, "\tilde\psi_i^{\omega}"] \arrow[dr, "\psi_i^{\omega}"'] & & \pi(P^{\omega}) \\
& \pi(Q_i) \arrow[ur, "(d_i^{\omega})^{-1}"'] &
\end{tikzcd}

\note{ce diagramme porte le numéro (38). À droite, une note rayée : « \ill{} \ill{} défini de l'isomorphisme $\Pi_{0,3} \simeq \Pi_{0,3}$ \ill{} », remplacée par « \uncertain{pour simplifier} ».}
Il suit que les relations (35) et (19) donnent
\[
(39)\quad \tilde\psi^{\omega}_{\omega i}(g) = \tilde\psi_i^{\omega}(\rho g) \quad \text{i.e.} \quad \tilde\psi_{\omega i} = \tilde\psi_i \circ \rho,
\]
\note{au-dessus de $\rho g$, entouré au crayon : « $\rho^{-1}$ ? ».}
et j'introduis les \add{deux} \uncertain{flèches}
\[
(40)\quad \tilde d_i^{\omega} : P^{\omega} \to P^{\omega'}, \qquad \tilde d_i^{\omega} = (d_i^{\omega'})^{-1} d_i^{\omega},
\]
permutés de façon simplement transitive par $\widetilde{\mathfrak{S}}_J$, \struck{et satisfaisant} satisfaisant
\[
(41)\quad \tau(\tilde d_i^{\omega}) = \tilde d_i^{\omega'}, \qquad \tilde d_i^{\omega} \tilde d_i^{\omega'} = 1
\]
On \struck{aura donc} \add{a} les relations
\[
(42)\quad \tilde d_i^{\omega}(\tilde\psi_i^{\omega}(g)) = \tilde\psi_i^{\omega'}(\sigma_\infty g).
\]
Revenant à l'\struck{automorphisme} \add{automorphisme} $u$, on aura \struck{\uncertain{automatiquement}}
\[
(43)\quad u(\tilde d_i^{\omega}) = \tilde\psi_i^{\omega'}(\sigma_\infty \tilde\beta)\, \tilde d_i^{\omega} = \tilde d_i^{\omega}\, \tilde\psi_i^{\omega}(\tilde\beta)
\]
avec
\[
(44)\quad \tilde\beta = \zeta\, \sigma_\infty(\zeta)^{-1}.
\]
\note{dans (43), le premier argument est surchargé ; on lit $\sigma_\infty\tilde\beta$.}
$\tilde\beta$ satisfaisant, \add{donc,} comme $\beta$ tantôt, la relation
\[
(45)\quad \tilde\beta\, \sigma_\infty(\tilde\beta) = 1
\]
Considérons maintenant le sous-groupoïde

\page{31}
\note{p.~13 de l'auteur}
\drawing{sphère avec les pôles $P_{+}$, $P_{-}$, les trois points marqués sur l'équateur et les méridiens $d_0^{\pm}$, $d_1^{\pm}$, $d_\infty^{\pm}$. Légende : $\mathrm{V}_J$.}
$\mathrm{IV}_J$, \add{induit sur les} \struck{formé des} deux sommets $P^{\omega}$ \struck{\ill{}} ($\omega \in \underline{\omega}$), on peut le regarder comme le groupoïde libre engendré par le graphe de \uncertain{noms} $\mathrm{V}_J$, ayant ces deux sommets, et trois arêtes — donc \uncertain{six} arcs $\tilde d_i^{\omega}$ — avec $\tilde d_i^{\omega}$ et $\tilde d_i^{\omega'}$ opposés — arcs permutés \struck{transitivement} \add{de façon} simpl\supplied{ement} tr\supplied{ansitive} \struck{\ill{}} par $\widetilde{\mathfrak{S}}_J$. Un \add{automorphisme} \struck{\ill{}} de ce groupoïde qui fixe les sommets est déterminé par les
\[
(46)\quad u(\tilde d_i^{\omega}) = \tilde d_i^{\omega} \tilde\psi_i^{\omega}(\tilde\beta)
\qquad \text{avec } \tilde\beta \in \widehat{\Pi}_{0,3} \text{ indépendant de } i, \omega
\]
la seule relation à satisfaire correspond aux équations $\tilde d_i^{\omega'} \tilde d_i^{\omega} = 1$, qui donne
\[
(47)\quad \tilde\beta\, \sigma_\infty(\tilde\beta) = 1
\]

Considérons enfin le groupoïde induit sur les \struck{pts} sommets $R_i^{\omega}$, $Q_i$, $P^{\omega}$ ($(6+3+2 = 11$ sommets$)$,

\page{32}
\note{p.~14 de l'auteur}
\drawing{sphère réunissant les figures précédentes : pôles $P^{+}$, $P^{-}$ ; sur l'équateur, $Q_0$, $Q_1$, $Q_\infty$ et, autour de chacun des trois points marqués, le petit cercle $R_i^{\pm}$ avec les arcs $a_i^{\pm}$ ; les arcs $c_i^{\pm}$ joignant les $Q$ aux $R$, et les méridiens $d_i^{\pm}$. Légende : $\mathrm{VI}$.}
c'est le groupoïde engendré par ces sommets, les arcs $a_i^{\omega}$, $c_i^{\omega}$, $d_i^{\omega}$ (trois paquets, \struck{\ill{}} chacun homogène sous $\widetilde{\mathfrak{S}}_J$) et les \add{six} relations
\[
(48)\quad (d^{\omega}_{\omega i})^{-1} (c^{\omega'}_{\omega i})^{-1} (a^{\omega}_{\omega i})^{-1} c_i^{\omega} d_i^{\omega} = 1
\]
\note{les indices de (48), surchargés, sont lus avec doute.}
L'\struck{automorphisme} \add{involution} $\tau$ est donnée par
\[
(49)\quad
\begin{cases}
\tau(R_i^{\omega}) = R_i^{\omega'}, \quad \tau Q_i = Q_i, \quad \tau(P^{\omega}) = P^{\omega'} \\
\tau(a_i^{\omega}) = (a_i^{\omega'})^{-1}, \quad \tau(c_i^{\omega}) = c_i^{\omega'}, \quad \tau(d_i^{\omega}) = d_i^{\omega'}
\end{cases}
\]
Les $\pi_1$ en les sommets $R_i^{\omega}$, $Q_i$, $P^{\omega}$ sont repérés respectivement, en termes de $\Pi_{0,3}$, par les isomorphismes
\[
(50)\quad
\begin{cases}
\Pi_{0,3} \xrightarrow{\ \varphi_i^{\omega}\ } \pi(R_i^{\omega}) \quad \text{\struck{$\to \pi($}} \\
\Pi_{0,3} \xrightarrow{\ \psi_i^{\omega}\ } \pi(Q_i) \\
\Pi_{0,3} \xrightarrow{\ \tilde\psi_i^{\omega}\ } \pi(P^{\omega})
\end{cases}
\]
reliés \struck{entre} \add{entre} eux par les relations suivantes, qui \struck{les} expriment les $\psi_i^{\omega}$, $\tilde\psi_i^{\omega}$ en termes des $\varphi_i^{\omega}$
\[
(51)\quad
\begin{cases}
d_i^{\omega}(\tilde\psi_i^{\omega}(g)) = \psi_i^{\omega}(g) \\
c_i^{\omega} \psi_i^{\omega}(g) = \varphi_i^{\omega}(g)
\end{cases}
\]
\note{devant chacune des deux lignes de (51), un symbole surchargé et rayé.}
On a, pour mémoire
\[
(52)\quad \psi_i^{\omega}(g) = \psi_i^{\omega'}(\sigma_\infty(g))
\]

\page{33}
\note{p.~15 de l'auteur}
\[
(53)\quad a_i^{\omega} \varphi_i^{\omega}(g) = \varphi_i^{\omega'}(\varepsilon_0(g))
\]
La donnée d'un endomorphisme des groupoïdes \add{profinis}, fixant les sommets et qui commute à $\widetilde{\mathfrak{S}}_J$, équivaut à celle des \struck{\ill{}} transformés des $a_i^{\omega}$, $c_i^{\omega}$, $d_i^{\omega}$, de façon compatible à l'action de $\widetilde{\mathfrak{S}}_J$, et de façon à satisfaire aux relations (48). On aura
\[
(54)\quad
\begin{cases}
u(a_i^{\omega}) = a_i^{\omega} \varphi_i^{\omega}(\lambda) \\
u(c_i^{\omega}) = \varphi^{\omega'}_{\omega' i}(\eta)\, c_i^{\omega} = c_i^{\omega} \psi_i^{\omega}(\eta) \\
u(d_i^{\omega}) = \psi_i^{\omega}(\zeta)\, d_i^{\omega} = d_i^{\omega} \tilde\psi_i^{\omega}(\zeta)
\end{cases}
\]
avec $\lambda, \zeta, \eta \in \widehat{\Pi}_{0,3}$, soumis à la seule relation déduite de (48). Au lieu de refaire le calcul, une fois de plus, pour expliciter la relation en $(\lambda, \eta, \zeta)$, il suffit \uncertain{de noter} que \struck{le} \add{le} groupoïde envisagé est la réunion des groupoïdes $\mathrm{II}$ et $\mathrm{IV}$, dont l'intersection est $\mathrm{III}$ \add{(ayant les sommets et commutant à $\mathfrak{S}_3$)}, donc \add{la donnée d'}un endomorphisme \add{équivaut à celle} d'un \add{tel} endom\supplied{orphisme} $(\lambda, \eta)$ de $\mathrm{II}$, et d'un endomorphisme $\zeta$ de $\mathrm{IV}$, de telle façon qu'ils se recollent \struck{\ill{}} $\mathrm{III}$ pour un $=$ endomorphisme $(\tilde\alpha)$ de $\mathrm{III}$, i.e.\ tels qu'on ait à la fois
\[
(55)\quad \tilde\alpha = \eta^{-1} \varepsilon_0(\lambda\eta) \quad \text{et} \quad \tilde\alpha = \zeta \rho^{-1}(\zeta^{-1})
\]

\page{34}
\note{p.~16 de l'auteur}
L'équation à satisfaire par $(\lambda, \eta)$ se réduit, \uncertain{comme} on \uncertain{sait}, à $\tilde\alpha \rho^{-1}(\tilde\alpha) \rho^{-2}(\tilde\alpha) = 1$ avec $\tilde\alpha$ donné en fonction de $(\lambda, \eta)$ par la première \struck{relation} \add{formule} (55) — mais cette relation \add{en $\tilde\alpha$} \struck{est} conséquence de la seconde relation (55), donnant $\tilde\alpha$ en termes de $\zeta$. Il reste donc, comme seule relation reliant $\lambda, \eta, \zeta$, que la relation (55) sous la forme (où ne figure plus $\tilde\alpha$)
\[
(56)\quad \eta^{-1} \varepsilon_0(\lambda\eta) = \zeta \rho^{-1}(\zeta)^{-1}
\]
Je vais la mettre sous une forme plus familière pour moi, en introduisant d'abord les quantités précédentes
\[
(57)\quad
\begin{cases}
\beta = \sigma_\infty(\eta)\eta^{-1} & (u(b_i^{\omega}) = b_i^{\omega} \varphi^{\omega}_{\omega i}(\beta)) \\
\alpha = \sigma_\infty(\lambda)\beta & (u(\tilde b_i^{\omega}) = \tilde b_i^{\omega} \varphi^{\omega}_{\omega i}(\alpha)) \\
\tilde\alpha = \eta^{-1}\varepsilon_0(\lambda\eta) = \zeta\rho^{-1}(\zeta)^{-1} & (u(\tilde a_i^{\omega}) = \psi^{\omega}_{\omega i}(\tilde\alpha)\, \tilde a_i^{\omega})
\end{cases}
\]
et en les modifiant ainsi, ainsi que $\zeta$ :
\[
(58)\quad
\begin{cases}
\beta' = \beta^{-1} \ (= \sigma_\infty(\beta)) = \eta\, \sigma_\infty(\eta)^{-1} \\
\alpha' = \alpha^{-1} = \beta^{-1} \sigma_\infty(\lambda)^{-1} = \beta' \sigma_\infty(\lambda^{-1})
\end{cases}
\qquad \text{donc } \beta' = \alpha' \sigma_\infty(\lambda)
\]
\struck{Je vais \ill{} modifier $\zeta$ en introduisant $\zeta' = $ \ill{} $\eta\,\zeta$ (\ill{} $\gamma$, \uncertain{grâce à} $\eta$) la relation (56) \uncertain{équivalant} \ill{}}
\note{deux lignes intercalées dans (58), puis rayées, avec un numéro (59) lui-même rayé.}
\struck{(59) Il \uncertain{reste} \uncertain{seulement} $\alpha'$, $\gamma$} décrits par les quantités $\boxed{\lambda, \eta, \zeta} \in \widehat{\Pi}_{0,3}$ (équivalentes à la donnée de $\lambda, \eta, \zeta$), \uncertain{auxquelles} j'associe $\alpha'$, $\beta'$ par
\[
(59)\quad
\begin{aligned}
\beta' &= \eta\, \sigma_\infty(\eta)^{-1} \quad \text{\struck{\ill{}}} \\
\alpha' &= \text{\struck{\ill{}}}\ \beta' \sigma_\infty(\lambda^{-1}) \quad \text{i.e.} \quad \beta' = \alpha' \sigma_\infty(\lambda)
\end{aligned}
\]

\page{35}
\note{p.~17 de l'auteur}
et ceci posé, la relation (56) équivaut à
\[
(60)\quad
\begin{cases}
\alpha' = \zeta' \rho^{-1}(\zeta')^{-1} \\
\text{i.e.}\quad \eta\, \sigma_\infty(\eta)^{-1} = (\zeta' \rho^{-1}(\zeta')^{-1})\, \sigma_\infty(\lambda)
\end{cases}
\]
\note{dans les deux lignes de (60), l'exposant de $\rho$ est surchargé ; on le lit $-1$.}
En effet, la formule (25\,bis) exprimant $\alpha$ en termes de $\tilde\alpha$
\[
\alpha = \rho(\eta \tilde\alpha \rho^{-1}(\eta)^{-1})
\]
\note{le renvoi est écrit « (75\,bis) » ; il s'agit de la formule (25\,bis) de la p.~26.}
se récrit, en passant aux inverses et posant
\[
\tilde\alpha' = \tilde\alpha^{-1}
\]
sous la forme
\[
(61)\quad \alpha' = \rho(\rho^{-1}(\eta)\, \tilde\alpha'\, \eta^{-1}) = \eta\, \rho(\tilde\alpha')\, \rho(\eta)^{-1}
\]
\struck{et la formule} (on considère ici $\alpha, \alpha', \tilde\alpha, \tilde\alpha'$ comme déduits des $\lambda, \eta \in \widehat{\Pi}_{0,3}$ par la première relation (55) pour ce qui est de $\tilde\alpha$, donc les autres se déduisent par (57), (61)). D'autre part, la \add{2e} formule (55) \struck{s'\ill{}} se récrit alors, en passant aux inverses
\[
\tilde\alpha' = \rho^{-1}(\zeta)\, \zeta^{-1}
\]
i.e.
\[
\rho(\tilde\alpha') = \zeta\, \rho(\zeta)^{-1}
\]
\struck{qui \ill{} équivaut (61) \ill{} s'écrit \ill{}}
\[
\alpha' = \eta\zeta\, \rho(\zeta^{-1}\eta^{-1}) \qquad \text{\struck{i.e.\ $\alpha'$ \ill{} (60)}}
\]
i.e.\ (60), avec
\[
\zeta' = \eta\zeta.
\]

\emph{En résumé} les automorphismes du groupoïde $\mathrm{VI}_J$, induisant l'identité sur les sommets et commutant à l'action de $\widetilde{\mathfrak{S}}_J$, correspondent aux
\[
(62)\quad \{(\lambda, \eta, \zeta') \in \widehat{\Pi}_{0,3} \times \widehat{\Pi}_{0,3} \times \widehat{\Pi}_{0,3}\}
\]
satisfaisant la relation des lacets :

\page{36}
\note{p.~18 de l'auteur}
\[
(63)\quad
\boxed{
\begin{aligned}
&\beta' = \alpha' \sigma_\infty(\lambda) \qquad (= \alpha' l_1^{\gamma} \ \text{si}\ \lambda = l_0^{\gamma}) \\
&\text{avec}\quad \alpha' = \zeta' \rho^{-1}(\zeta')^{-1} \\
&\phantom{\text{avec}\quad} \beta' = \eta\, \sigma_\infty(\eta)^{-1}
\end{aligned}
}
\]
\note{dans la deuxième ligne de (63), l'exposant de $\rho$ est surchargé ; on le lit $-1$ comme en (60).}
\marginal{(en biais, à gauche) le passage de $u$ à $\tau(u)$ \uncertain{s'écrit} : $\lambda \mapsto \varepsilon_0^{-1}(\tau_\infty(\lambda^{-1}))$ ($\lambda \mapsto \lambda$ si $\lambda = l_0^{\gamma}$), $\eta \mapsto \tau_\infty(\eta)$, $\zeta \mapsto \sigma_\infty\tau_\infty(\zeta) = \tilde\sigma_\infty(\zeta)$}
\marginal{(en biais, à droite) NB. L'involution \uncertain{\ill{}} \ill{} \ill{} : $\zeta' \mapsto \rho H \zeta'$, $\eta \mapsto l_0 H \eta$, i.e.\ $\alpha' \mapsto l_0 H \alpha'$ \ill{}, $\beta' \mapsto l_0 H \beta'$ \ill{}}
\note{les deux notes marginales en biais de cette page se lisent mal ; la seconde, en particulier, est très incertaine.}

Ici, les quantités $\alpha', \beta', \zeta'$ ont été introduites en termes des quantités primitives $\lambda, \eta, \zeta$ ($\mapsto \beta, \alpha, \tilde\alpha$), s'introduisant dans les formules (54) (et (57)), par le pur calcul \struck{\ill{}} \add{dans} $\widehat{\Pi}_{0,3}$, comme donnant lieu aux formules les plus sympathiques. Mais on peut les interpréter aussi de façon « géométrique », en posant \struck{\ill{}} \uncertain{comme} on \uncertain{a fait} \uncertain{pour} $\beta', \alpha', \tilde\alpha'$
\[
(64)\quad b_i'^{\omega} = (b_i^{\omega})^{-1} \ (= b_i^{\omega'}), \qquad \tilde b_i'^{\omega} = (\tilde b_i^{\omega})^{-1}, \qquad \tilde a_i'^{\omega} = (\tilde a_i^{\omega})^{-1}
\]
de sorte que les formules (57) deviennent
\[
(65)\quad
\begin{cases}
u(b_i'^{\omega}) = \varphi^{\omega}_{\omega i}(\beta')\, b_i'^{\omega} \\
u(\tilde b_i'^{\omega}) = \varphi^{\omega}_{\omega i}(\alpha')\, \tilde b_i'^{\omega} \\
u(\tilde a_i'^{\omega}) = \tilde a_i'^{\omega}\, \psi^{\omega}_{\omega i}(\tilde\alpha')
\end{cases}
\]
Enfin, la quantité $\zeta' = \eta\zeta$ est obtenue géométriquement, en introduisant les chemins
\[
(66)\quad \tilde d_i'^{\omega} : P^{\omega} \to R_i^{\omega'}, \qquad \tilde d_i'^{\omega} = c^{\omega}_{\omega i} \cdot d^{\omega}_{\omega i}
\]
\note{sous la flèche de (66), le chemin est détaillé : $P^{\omega} \xrightarrow{d^{\omega}_{\omega i}} Q_{\omega i} \xrightarrow{c^{\omega}_{\omega i}} R_i^{\omega'}$.}
\drawing{en bas à gauche : petit triangle $P^{\omega}$, $Q_{\omega i}$, $R_i^{\omega'}$ ; arcs $d^{\omega}_{\omega i}$ de $P^{\omega}$ à $Q_{\omega i}$, $c^{\omega}_{\omega i}$ de $Q_{\omega i}$ à $R_i^{\omega'}$, et le chemin composé $\tilde d_i'^{\omega}$ ; une boucle autour de $R_i^{\omega'}$.}
par la formule

\page{37}
\note{p.~19 de l'auteur}
\[
(67)\quad u(\tilde d_i'^{\omega}) = \varphi_i^{\omega'}(\underbrace{\zeta'}_{\eta\zeta})\, \tilde d_i'^{\omega}
\]
On a aussi des chemins
\[
(66')\quad d_i^{\omega} : P^{\omega} \to R_i^{\omega}, \qquad d_i^{\omega} = c^{\omega'}_{\omega' i} \cdot d^{\omega}_{\omega' i}
\]
\note{sous la flèche de (66$'$) : $P^{\omega} \xrightarrow{d^{\omega}_{\omega' i}} Q_{\omega' i} \xrightarrow{c^{\omega'}_{\omega' i}} R_i^{\omega}$. Le nom $d_i^{\omega}$ est celui qu'il écrit, bien qu'il ait déjà servi pour les arcs de $P^{\omega}$ à $Q_i$.}
\drawing{à gauche : $P^{\omega}$ en haut, $Q_{\omega' i}$ en bas, $R_i^{\omega}$ à gauche ; flèche $d^{\omega}_{\omega' i}$ de $P^{\omega}$ à $Q_{\omega' i}$, arc $c^{\omega'}_{\omega' i}$ de $Q_{\omega' i}$ à $R_i^{\omega}$, et en pointillé le chemin composé.}
\marginal{NB $\tau(d_i^{\omega}) = d_i'^{\omega}$}
et on aura également
\[
(67')\quad u(\tilde d_i^{\omega}) = \varphi_i^{\omega}(\eta\, \sigma_\infty(\zeta))\, \tilde d_i^{\omega}
\]
\note{l'argument de (67$'$) est surchargé ; la lecture $\eta\,\sigma_\infty(\zeta)$ est incertaine.}

\drawing{sphère : les deux pôles, les huit sommets $R_i^{\omega}$, $P^{\omega}$, et les arcs $d_i^{\omega}$, $d_i'^{\omega}$ (méridiens, pleins ou pointillés), $a_i^{\omega}$ (petits cercles) et $b_i^{\omega}$ (le long de l'équateur). Légende : $\mathrm{VII}$.}
\marginal{NB les $b_i^{\omega}$ reliant directement les $R$ \uncertain{les} $Q_i$ sont ici inutiles}
Je vois encore deux cas à regarder, savoir les sous-groupoïdes pleins de $\mathrm{VI}$ induits sur les \add{huit} sommets $(R_i^{\omega}, P^{\omega})$, soit $\mathrm{VII}$, décrit par le graphe de \uncertain{noms} et les flèches $d_i'^{\omega}$, $d_i^{\omega}$, $a_i^{\omega}$, $b_i^{\omega}$ (24 flèches), et celui induit sur les \struck{\ill{} points} \add{sommets} $Q_i$, $P^{\omega}$ — à strictement parler, c'est le groupoïde déjà nommé $\mathrm{IV}$ \uncertain{ayant} les \uncertain{mêmes} \struck{\ill{}} sommets, mais représenté par un graphe un peu plus grand — dans \uncertain{$\mathrm{VIII}$} on distingue les \struck{flèches} \add{\uncertain{six}} $d_i^{\omega}$ seulement, et on ne donne pas en plus les $\tilde a_i^{\omega}$, ce qui est inutile en effet par (35). (Dans il \ill{} \uncertain{tout} donc (37)), \struck{\ill{}} la donnée de $\tilde\alpha$, par
\drawing{seconde sphère, au crayon : les deux pôles reliés par trois méridiens, les trois points marqués sur l'équateur. Légende : $\mathrm{VIII} = \mathrm{IV}$.}
\note{la phrase se poursuit à la p.~38.}

\page{38}
\note{p.~20 de l'auteur}
\add{les} repère \add{les} $u(\tilde a_i^{\omega})$, \struck{résulte} déjà de celle de $\zeta$, qui repère \add{les} $u(d_i^{\omega})$. // Les graphes $\mathrm{I}$ à $\mathrm{VII}$ donnent lieu au diagramme d'inclusions

\begin{tikzcd}[column sep=small, row sep=small]
& & \mathrm{III} \arrow[dl, hook'] \arrow[dr, hook] & & \\
\mathrm{I} \arrow[r, hook] \arrow[dr, hook] & \mathrm{II} \arrow[dr, hook] & & \mathrm{IV} \arrow[dl, hook'] & \mathrm{V} \arrow[l, hook'] \\
& \mathrm{VII} \arrow[r, hook] & \mathrm{VI} & &
\end{tikzcd}

\note{diagramme (68). Sur la page, $\mathrm{I}$, $\mathrm{II}$, $\mathrm{VII}$, $\mathrm{VI}$ sont encadrés ensemble ; une flèche partant de $\mathrm{II}$ vers la droite est rayée.}
Par le calcul, je peux déjà prédire que la donnée d'un automorphisme de $\mathrm{VII}$ fixant les sommets et commutant à $\widetilde{\mathfrak{S}}_J$, équivaut à la donnée des quantités
\[
(69)\quad (\lambda, \beta', \zeta') \in \widehat{\Pi}_{0,3} \times \widehat{\Pi}_{0,3} \times \widehat{\Pi}_{0,3}
\]
\marginal{donc $\beta'$ inutile !}
reliées par les équations \struck{(63) i.e.}
\[
(70)\quad
\begin{cases}
\beta' = \alpha' \sigma_\infty(\lambda) \quad \text{où} \quad \alpha' = \zeta' \rho^{-1}(\zeta')^{-1} \\
\beta' \sigma_\infty(\beta') = 1 \quad \text{i.e.} \quad \alpha' \sigma_\infty(\lambda)\, \sigma_\infty(\alpha')\, \lambda = 1
\end{cases}
\]
\note{la seconde ligne de (70) est ajoutée en interligne ; son dernier facteur est surchargé.}
(tandis que \add{dans le cas} \struck{\ill{}} $\mathrm{II}$ \add{$u$} est défini par \struck{$(\lambda, \eta, \alpha')$},
\[
(69')\quad (\lambda, \eta, \alpha') \in \widehat{\Pi}_{0,3} \times \widehat{\Pi}_{0,3} \times \widehat{\Pi}_{0,3}
\]
\struck{reliés par} avec les équations
\[
(70')\quad
\begin{cases}
\beta' = \alpha' \sigma_\infty(\lambda) \quad \text{où} \quad \beta' = \eta\, \sigma_\infty(\eta)^{-1} \\
\alpha' \rho(\alpha') \rho^{2}(\alpha') = 1
\end{cases}
\]
\marginal{donc $\alpha'$ inutile !}
\note{dans (69$'$), $\lambda$ et $\eta$ sont soulignés ; la parenthèse ouverte avant « tandis que » n'est pas refermée sur la page.}

\page{39}
\note{p.~21 de l'auteur}
Je vais expliciter les calculs \add{\uncertain{faits}} dans le groupoïde $\mathrm{VII}$, associé au graphe de \uncertain{noms}, \ill{} des flèches génératrices
\[
(71)\quad
\begin{cases}
\tilde d_i^{\omega} : P^{\omega} \to R_i^{\omega} \\
\tilde d_i'^{\omega} : P^{\omega} \to R_i^{\omega'} \\
a_i^{\omega} : R_i^{\omega} \to R_i^{\omega'}
\end{cases}
\]
liées par les relations
\[
(72)\quad
\begin{cases}
(\tilde d_i'^{\omega})^{-1} a_i^{\omega} \tilde d_i^{\omega} = 1 \\
(\tilde d'^{\omega}_{\omega i})^{-1} \tilde d^{\omega'}_{\omega i} (\tilde d_i'^{\omega'})^{-1} \tilde d_i^{\omega} = 1
\end{cases}
\]
\note{les indices et exposants de la seconde ligne de (72) sont surchargés et lus avec doute.}
\marginal{six relations indépendantes exprimant les $\tilde d_i'^{\omega}$ en termes des $\tilde d_i^{\omega}$ et $a_i^{\omega}$, ou les $a_i^{\omega}$ en termes des $\tilde d_i^{\omega}, \tilde d_i'^{\omega}$. \uncertain{En tout} six relations, mais il n'y en a que trois indépendantes, car l'équation pour $i, \omega$ et pour $(\omega i, \omega')$ est la même.}
\note{cette note, à droite de (72), se rapporte à la seconde ligne ; sa première partie est insérée en interligne.}
\drawing{sphère : pôles $P^{+}$, $P^{-}$ ; autour des trois points marqués, les petits cercles $R_i^{\pm}$ avec les arcs $a_i^{\pm}$ ; les méridiens $\tilde d_i^{\pm}$ et $\tilde d_i'^{\pm}$ ($i = 0, 1, \infty$) descendant des pôles vers les $R$, en trait plein ou pointillé selon qu'ils sont vus ou cachés.}
On les repère \uncertain{ensuite} par les équations
\[
(73)\quad
\begin{cases}
u(\tilde d_i^{\omega}) = \varphi_i^{\omega}(\zeta_1)\, \tilde d_i^{\omega} \\
u(\tilde d_i'^{\omega}) = \varphi_i^{\omega'}(\zeta')\, \tilde d_i'^{\omega} \\
u(a_i^{\omega}) = a_i^{\omega} \varphi_i^{\omega}(\lambda)
\end{cases}
\qquad \zeta', \zeta_1, \lambda \in \widehat{\Pi}_{0,3}
\]
\note{dans les deux premières lignes de (73), un premier facteur, raturé à l'encre, est remplacé par celui qui est donné ici.}
et il reste à exprimer que $u$ satisfait les relations (72). La première s'explicite sous la forme
\[
(74)\quad \zeta' = \varepsilon_0(\lambda\zeta_1) \qquad \text{i.e.} \qquad \zeta_1 = \lambda^{-1} \cdot \varepsilon_0^{-1}(\zeta')
\]
et on trouve également
\[
(75)\quad \beta\, \zeta_1 = \beta' \sigma_\infty(\zeta')
\]
où $\beta'$ défini par la condition (65\uncertain{a})\note{le premier $\beta$ de (75) est surchargé.}

Cette relation permettant de tirer $\beta'$ en termes
\note{la phrase se poursuit à la p.~40.}

\page{40}
\note{p.~22 de l'auteur}
de $\zeta', \zeta_1$ — ou de \uncertain{l'une} des \uncertain{deux} quantités seulement
\[
(75)\quad
\begin{aligned}
\beta' &= \zeta_1 \sigma_\infty(\zeta')^{-1} = \text{\struck{$\varepsilon_0(\lambda\zeta_1)$}}\ \zeta_1 (\sigma_\infty\varepsilon_0)(\lambda\zeta_1)^{-1} \\
&= \lambda^{-1} \cdot \varepsilon_0^{-1}(\zeta')\, \sigma_\infty(\zeta')^{-1}.
\end{aligned}
\]
\struck{\ill{} la dernière} \add{ces} expressions de $\beta'$ \uncertain{on tire}
\[
(75')\quad \sigma_\infty(\beta') = \sigma_\infty(\zeta_1)\, \zeta'^{-1} = \sigma_\infty(\zeta_1)\, \varepsilon_0(\lambda\zeta_1)^{-1}
= \sigma_\infty(\lambda)^{-1} \overbrace{\rho(\zeta')\, \zeta'^{-1}}^{\alpha'^{-1}}
\]
\note{le numéro (75) sert deux fois, ici et à la p.~39 ; à la fin de la ligne de (75$'$), un groupe raturé.}
et la troisième expression de $\beta'$ donne \struck{\ill{}}
\struck{(77) Ce dit, l'équation pour (72\,b) équivalente :}
\[
\text{\struck{$\beta' \sigma_\infty(\beta') = 1$}}
\]
\struck{\ill{} $\sigma_\infty(\beta')$ en posant $\beta_1 = \sigma_\infty(\beta')$ \ill{}}
\note{ce passage est encadré puis rayé ; des traits ondulés au crayon le traversent.}
s'écrit aussi
\[
(77)\quad \sigma_\infty(\beta')^{-1} = \alpha' \sigma_\infty(\lambda).
\]
Ceci posé, l'équation (72\,b) équivaut à
\[
(78)\quad \beta' \sigma_\infty(\beta') = 1 \qquad \text{ou encore} \quad \sigma_\infty(\beta')^{-1} = \beta'
\]
\note{le numéro (78) est surchargé.}
ce qui est invariant par la substitution $\beta' \mapsto \sigma_\infty(\beta')$ \struck{dans} $\beta' \mapsto \beta'^{-1}$, donc posant (dans ses \uncertain{simplifications} \uncertain{($\ast$)})
\[
\beta_1 = \sigma_\infty(\beta')^{-1} = \alpha' \sigma_\infty(\lambda)
\]
l'équation (78) équivaut à l'équation
\[
(\ast\ast)\quad \beta_1 \sigma_\infty(\beta_1) = 1
\]
et implique que $\beta_1 = \sigma_\infty(\beta_1)^{-1}$ donc $\beta_1 = \beta'$. En résumé, on peut donc se \uncertain{servir} \uncertain{plutôt} de
\[
(78)\quad \beta' = \alpha' \sigma_\infty(\lambda) \qquad \text{où} \quad \alpha' = \zeta' \rho(\zeta')^{-1}
\]
et la relation à imposer : $(\lambda, \zeta')$ est, \uncertain{au choix},
\note{le numéro (78) sert à nouveau, surchargé. Le texte continue au-delà de la p.~40, hors de ce lot.}
\end{document}
